\section{预训练语料与样本准备}
\label{sec:nlp-llm-data}

概率模型拟合的是训练过程中实际看到的序列分布。网页数量很多，不等于模型得到的语言、事实和推理材料就均衡；同一篇通知被转载一百次，也没有增加一百种独立的事件。准备语料的任务，是从原始材料中形成有边界、可追踪且能够按指定比例取用的训练样本。模型计算尚未开始，学习对象已经在这些选择中发生变化。

\subsection{来源选择与能力覆盖}
\label{sec:nlp-llm-data-sources}

设候选材料来自网页、书籍、代码和专业文档等来源。网页能够包含日常表达和近期信息，代码能够包含形式化约束与程序结构，专业文档能够包含通用网页较少涉及的概念。这里的“来源”描述材料怎样产生，并不自动给出质量或能力标签：网页可能是完整教程，也可能只有导航菜单；代码中也可能有过时接口或错误示例。来源标签的作用，是让后续处理能够检查某种材料是否被过度删除或过度重复。

为了把“覆盖”说清楚，可以为每篇文档保留语言、时间、来源、主题和长度等属性，再分别统计它们的分布。若希望模型处理中文活动通知和英文技术文档，至少需要检查两类材料是否存在、包含哪些表达以及是否覆盖相关时间范围。材料中有某个主题，仍不意味着模型能正确完成该主题的任务，因为学习还取决于采样次数、目标和参数容量。覆盖检查排除的是明显的信息缺口，不是能力保证。

材料的时间边界尤其影响事实性判断。假设训练语料中只有“读书会于周四举行”，而新的更改通知没有进入语料，降低旧材料上的预测损失不能使模型知道周五的新安排。这类缺口需要新材料或者使用时提供外部证据，不能只通过增大网络消除。为比较不同模型，应记录其材料时间范围，避免把接触过更新信息误判为相同信息条件下的建模优势。

来源选择还需要确定哪些材料可以进入当前训练集，以及敏感内容和个人信息如何处理。这些约束改变可使用的文档集合，因而属于数据准备本身。保留来源记录有助于发现问题后追溯处理规则；后续安全示范与偏好学习另有职责，不能替代这里的材料检查。

\subsection{质量筛选与重复控制}
\label{sec:nlp-llm-data-filtering}

候选文档中常混有乱码、模板、极短片段和重复正文。若直接按文档大小取样，模型会花费大量更新去拟合这些重复或低信息量的模式。质量筛选先依据明确规则或评分保留材料；重复控制再限制相同内容在训练分布中的权重。两项操作都改变训练对象，不能仅被看成不影响学习的文件清理。

规则筛选可以检查文本可解析性、语言比例、重复行或异常字符。模型评分则先定义希望保留的特征，例如文档是否连贯、是否包含解释性内容，再用分类器或语言模型估计分数。评分器的判断也有误差：只偏好规范书面语，可能删除口语、方言和少见专业表达。阈值应通过被保留与被删除样本的抽查来解释，不能把较高分数直接当作真实知识更正确。

完全重复可以通过规范化后的文本哈希识别。近似重复则需要比较文档或片段的重合，例如把文档表示为连续词元片段集合，再依据集合重合率判定是否属于同一重复簇。去重粒度会改变结果：只删完全相同的整篇文档，会留下含不同页眉的转载；激进删除共享段落，又可能误伤不同文档都需要的定义或代码模板。训练集与评价集之间也要检查重合，但这种检查不能证明没有语义改写后的污染。

DataComp-LM将数据筛选放在共同的模型、预算和评价流程下比较，FineWeb则公开分析了过滤与去重选择。它们提供的认识是，数据处理应通过可比较的学习结果检验，而不能只报告删掉多少文本；其具体收益仍受语言范围、模型规模和评价任务约束\scite{nlp-llm-dclm2024,nlp-llm-fineweb2024}。在本章的活动通知例子中，保留一份正确正文、删除九份相同转载，会降低该通知的重复曝光，但不会自动增加模型对其他活动安排的认识。

\subsection{词元化与文档序列组织}
\label{sec:nlp-llm-tokenization}

筛选后的字符串还不是网络使用的整数序列。\term{词元}（token）是模型一次读取或预测的离散单位；\term{词元化}（tokenization）把文本转换为这些单位的序列。词元可以对应字、常见子词或字节片段。这里需要区分文字、语言学意义上的词和模型单位：“下午两点”可以被分成多个词元，也可以部分合并，具体取决于词元化规则。

设编码函数将文档$d$映射为$x_{1:T}=(x_1,\ldots,x_T)$，其中$x_t\in\mathcal V$，$\mathcal V$是有限词表。输入嵌入和输出分布都以这个词表为接口。采用子词单位时，常见片段可以共用一个编号，稀有字符串再拆成较短单位；采用字节单位时，基本字节集合固定，但同一段文本往往需要更多步。一个词元覆盖多少字符，影响序列长度和可容纳的材料，因而不能直接用不同词元化器下的每词元损失比较文本建模质量。

文档组织还要说明模型在哪里开始、在哪里结束，以及相邻文档能否彼此读取。可以在通知结尾加入结束标记，将多个短文档装入一个固定长度序列；也可以为各文档设置相互隔离的注意掩码。前者若允许跨文档读取，就改变了条件信息：后一个文档的预测能够利用前一个文档。结束标记表达边界，但本身不等于强制隔离。切分长文档时，也应检查需要跨段的信息是否仍能共同出现。

例如，假设“读书会／改到／周五／下午／两点／举行”是一次教学词元化，最大长度只有四。若永远切成前四个与后两个词元，“两点”所在样本可能失去“改到周五”的条件。重叠切分可以保留部分上下文，但会增加某些词元的重复计分。具体监督位置由\ref{sec:nlp-llm-objective-positions}小节确定；在这里先保留切分与文档边界，才能知道后续目标用了什么信息。

\subsection{多源混合与采样配比}
\label{sec:nlp-llm-data-mixture}

完成统一编码后，材料的原始数量仍不等于训练时的出现比例。设处理后的$K$类来源分别给出样本分布$P_1,\ldots,P_K$，训练时先以概率$\alpha_k$选择来源$k$，再从该来源取样。实际训练分布为
\begin{equation}
 P_{\bm{\alpha}}(x)=\sum_{k=1}^{K}\alpha_k P_k(x),
 \qquad \alpha_k\geq 0,\quad\sum_{k=1}^{K}\alpha_k=1.
 \label{eq:nlp-llm-mixture}
\end{equation}
权重$\alpha_k$描述训练曝光，而$P_k$包含该来源内部的筛选、切分和取样方式。若按词元预算分配混合，权重应按所消耗的词元计量；若按样本抽取而样本长度不同，最终词元比例通常不会等于样本比例。

假设通用文本有九十万词元，专业文档有十万词元。按原始数量抽取，专业材料约占一成；若训练词元各分一半，专业材料会被重复使用更多次。提高专业权重既增加相关模式的学习机会，也减少同一预算内通用材料的曝光，并可能增加小语料的记忆。这里没有仅由材料数量决定的最优比例，需要说明希望改善哪些评价对象。

若为每种候选比例都完整训练大模型，比较成本很高。DoReMi先在较小代理模型上，根据各来源相对于参考模型的额外损失学习重加权，再将得到的比例用于较大模型；参考损失有助于避免单纯追逐本来就更难预测的来源。RegMix改为在许多比例上训练小模型，拟合“比例到评价结果”的回归关系，再选择预测较好的比例\scite{nlp-llm-doremi2023,nlp-llm-regmix2024}。前者利用代理训练中的损失反馈，后者利用跨候选配置的预测，两者都需要检验小规模判断能否迁移到目标训练。

混合比例也可以随阶段变化，例如在较广的材料上建立通用参数后，再提高某个领域的占比。只要比例或来源发生变化，就应重新说明训练分布，而不能把整个过程概括成使用了某个固定语料集。现在我们得到了可取样的词元序列，下一节讨论模型怎样把这些序列变成条件概率。
